%  LaTeX support: latex@mdpi.com 
%  For support, please attach all files needed for compiling as well as the log file, and specify your operating system, LaTeX version, and LaTeX editor.

%=================================================================
\documentclass[applsci,article,submit,pdftex,moreauthors]{Definitions/mdpi} 

%=================================================================
% MDPI internal commands
\firstpage{1} 
\makeatletter 
\setcounter{page}{\@firstpage} 
\makeatother
\pubvolume{1}
\issuenum{1}
\articlenumber{0}
\pubyear{2022}
\copyrightyear{2022}
%\externaleditor{Academic Editor: Firstname Lastname}
\datereceived{} 
%\daterevised{} % Only for the journal Acoustics
\dateaccepted{} 
\datepublished{} 
%\datecorrected{} % Corrected papers include a "Corrected: XXX" date in the original paper.
%\dateretracted{} % Corrected papers include a "Retracted: XXX" date in the original paper.
\hreflink{https://doi.org/} % If needed use \linebreak
%\doinum{}
%------------------------------------------------------------------
% The following line should be uncommented if the LaTeX file is uploaded to arXiv.org
%\pdfoutput=1

%=================================================================
% Add packages and commands here. The following packages are loaded in our class file: fontenc, inputenc, calc, indentfirst, fancyhdr, graphicx, epstopdf, lastpage, ifthen, lineno, float, amsmath, setspace, enumitem, mathpazo, booktabs, titlesec, etoolbox, tabto, xcolor, soul, multirow, microtype, tikz, totcount, changepage, attrib, upgreek, cleveref, amsthm, hyphenat, natbib, hyperref, footmisc, url, geometry, newfloat, caption

\graphicspath{{figures/}} % path to folder with figures
\usepackage{caption}
\usepackage{subcaption}
%\usepackage[export]{adjustbox}
\usepackage{listings}
\usepackage{xcolor}
\usepackage{graphicx}
\usepackage{gensymb} % for degree symbol
\usepackage{lscape}
\usepackage{pxfonts}
\usepackage{graphicx}
\usepackage{tikz}
	\newlength\imagewidth
	\newlength\imagescale
\usepackage{tabularx}
\usepackage[autostyle]{csquotes}
\newcolumntype{Y}{>{\centering\arraybackslash}X}
\usepackage{booktabs,caption}
\usepackage[flushleft]{threeparttable}

\definecolor{codegreen}{rgb}{0,0.6,0}
\definecolor{codegray}{rgb}{0.5,0.5,0.5}
\definecolor{codepurple}{rgb}{0.58,0,0.82}
\definecolor{backcolour}{RGB}{255,255,255}
\definecolor{SlateGray}{HTML}{708090}

\lstdefinestyle{mystyle}{
    backgroundcolor=\color{backcolour},   
    keywordstyle=\bfseries,
    numberstyle=\tiny\color{codegray},
    stringstyle=\color{SlateGray},
    basicstyle=\ttfamily\footnotesize,
    breakatwhitespace=false,         
    breaklines=true,                 
    captionpos=t, 
    keepspaces=true,                 
    numbers=left,                    
    numbersep=5pt,                  
    showspaces=false,                
    showstringspaces=false,
    showtabs=false,                  
    tabsize=2,
    frame=topline|bottomline,
    }

\lstset{style=mystyle}
\captionsetup[lstlisting]{position=top, margin=0cm,labelfont={bf, small, stretch=1.17}, labelsep=colon, textfont={small, stretch=1.17}, aboveskip=6pt, singlelinecheck=off, justification=justified}
	
\usepackage{soul} % highlighting text

%=================================================================
%% Please use the following mathematics environments: Theorem, Lemma, Corollary, Proposition, Characterization, Property, Problem, Example, ExamplesandDefinitions, Hypothesis, Remark, Definition, Notation, Assumption
%% For proofs, please use the proof environment (the amsthm package is loaded by the MDPI class).

%=================================================================
% Full title of the paper (Capitalized)
\Title{R Libraries for Remote Sensing Data Classification by k-means Clustering and NDVI Computation in Congo River Basin, DRC}

% MDPI internal command: Title for citation in the left column
\TitleCitation{R Libraries for Remote Sensing Data Classification by k-means Clustering and NDVI Computation in Congo River Basin, DRC}

% Author Orchid ID: enter ID or remove command
\newcommand{\orcidauthorA}{0000-0002-5759-1089} % Add \orcidA{} behind the author's name
\newcommand{\orcidauthorB}{0000-0002-6461-1551} % Add \orcidB{} behind the author's name

% Authors, for the paper (add full first names)
\Author{Polina Lemenkova $^{1}$\orcidA{} and Olivier Debeir $^{1}\orcidB{}$}
% \dagger,\ddagger
%\longauthorlist{yes}

% MDPI internal command: Authors, for metadata in PDF
\AuthorNames{Polina Lemenkova $^{1}$ and Olivier Debeir $^{1}$}

% MDPI internal command: Authors, for citation in the left column
\AuthorCitation{Lemenkova, P.; Debeir, O.}
% If this is a Chicago style journal: Lastname, Firstname, Firstname Lastname, and Firstname Lastname.

% Affiliations / Addresses (Add [1] after \address if there is only one affiliation.)
\address{%
$^{1}$ \quad Université Libre de Bruxelles, École polytechnique de Bruxelles (EPB, Brussels Faculty of Engineering), Laboratory of Image Synthesis and Analysis, Building L, Campus de Solbosch, ULB -- LISA CP165/57, Avenue Franklin D. Roosevelt 50, B-1050, Brussels, Belgium}

% Contact information of the corresponding author
\corres{Correspondence: polina.lemenkova@ulb.be; Tel.: +32471860459 (P.L.)}

% Abstract (Do not insert blank lines, i.e. \\) 
\abstract{In this paper, an image analysis framework is formulated for Landsat-8 OLI-TIRS scenes using R programming language. The libraries of R were shown to be effective in remote sensing data processing tasks such as classification using k-means clustering and computing the Normalized Difference Vegetation Index (NDVI). The data handling is processed using the integration of RStoolbox, terra, raster, rgdal and auxiliary packages of R. The proposed approach of image processing using R is designed to exploit the parameters of image bands as cues to detect land cover types and vegetation parameters corresponding to the spectral reflectance of the objects represented on the Earth's surface. Our method is effective for processing the time-series of the images taken in various periods to monitor the landscape dynamics in middle part of Congo River basin, D.R.C. Whereas previous approaches primarily used GIS software, we proposed to explicitly use the scripting methods for satellite image analysis by applying the extended functionality of R. The application of scripts for geospatial data is effective and robust method contrasting to the traditional approaches due to the high automation and machine-based graphical processing. The algorithms of R libraries are adjusted to spatial operations such as projection and transformation, object topology, classification and map algebra. The data included Landsat-8 OLI-TIRS covering the three regions along the river of Congo: Bumba, Basoko and Kisangani for the years 2013, 2015 and 2022. We also validated the performance of graphical data handling for cartographic visualization using R libraries for visualising changes in land cover types by k-means clustering and calculation of the NDVI for vegetation analysis.
}
% Keywords
\keyword{image processing; remote sensing; Landsat; R language; programming; cartography; mapping; data visualization; NDVI ; Africa} 

% The fields PACS, MSC, and JEL may be left empty or commented out if not applicable
\PACS{91.10.Da; 91.10.Jf; 91.10.Sp; 91.10.Xa; 96.25.Vt; 91.10.Fc; 95.40.+s; 95.75.Qr; 95.75.Rs; 42.68.Wt}
\MSC{86A30; 86-08; 86A99; 86A04}
\JEL{Y91; Q20; Q24; Q23; Q3; Q01; R11; O44; O13; Q5; Q51; Q55; N57; C6; C61}

%%%%%%%%%%%%%%%%%%%%%%%%%%%%%%%%%%%%%%%%%%
% Only for the journal Diversity
%\LSID{\url{http://}}

%%%%%%%%%%%%%%%%%%%%%%%%%%%%%%%%%%%%%%%%%%
% Only for the journal Applied Sciences
\featuredapplication{Libraries of R programming language \emph{RStoolbox}, \emph{raster} and \emph{terra} were used to classify Landsat OLI-TIRS satellite images by \emph{k}-means clustering and NDVI computation for comparative vegetation mapping of tropical rainforests for the period of 2013-2022 in the three target areas (Bumba, Basoko, Kisangani) of middle Congo River Basin, central Africa, D.R.C.}
%%%%%%%%%%%%%%%%%%%%%%%%%%%%%%%%%%%%%%%%%%

\begin{document}

%%%%%%%%%%%%%%%%%%%%%%%%%%%%%%%%%%%%%%%%%%
\section{Introduction}
Satellite image processing is a fundamental problem in remote sensing data analysis and is a key step towards the automatic visualization of the Earth's landscapes. The significance of the satellite images for visualising the Earth is that they incorporate rich information about the land cover types which can be derived from interpretation of the pixels. The brightness of pixels in the satellite images correspond to the spectral reflectance of various surfaces and objects: built-up urban areas, bare soil, deserts, forest and various vegetation types, water areas, etc. Since the absorption of solar radiation varies in these objects, their spectral reflection is different. Such fundamental properties of optics of the materials form a basis for the environmental applications of satellite imagery generated by the sensors onboard of the satellites \cite{Lillesand}. The differences in spectral signatures of the surfaces are detectable in bands, or channels, of the satellite imagery and digital numbers (DNs) of the pixels in bitmap images which can be used for land cover mapping \cite{Manakos}. 

Satellite images processed by remote sensing tools, are often used to visualise landscape dynamics as a way to detect land cover changes through exploiting the similarities and differences in pixels of the images in the Landsat \cite{SIKUZANI2020e01333,DUVEILLER20081969} or Sentinel \cite{9554698,Lemenkova202021,9553905} or their integration \cite{ZHANG2021112470,CHEN2021102386}. Some of the methods of image processing use time series analysis using fine spatial resolution imagery \cite{7326202,6352722}, topological data extraction by spectral analysis derived from SAR \cite{1027205,520569},\hl{ }photogrammetry\hl{ or }LiDAR processing to analyse forest canopy \hl{by collecting 3D data }\cite{rs13183777}. Additionally, image processing may use supervised classification approach which includes automated extraction of the training samples and reference maps \cite{rs12172705}, regression tree analysis for canopy height modelling \cite{HANSEN2016221}, or statistical methods to analyse forest degradation \cite{7326141}.

Spectral bands that comprise the satellite imagery consist of a matrix of pixels. Different intensity of the wavelengths of light of the pixels is visible on the image as colours of various hues and shadowing. The Landsat OLI/TIRS consists of 11 bands including spectral bands: ultra violet (UV), visual, near-infrared (NIR), short-wave infrared (SWIR) and thermal infrared (TIR). Red and NIR bands, which correspond to channels 4 and 5, are specifically useful for vegetation analysis and forest mapping as has been demonstrated in related works \cite{DALPONTE2020102206,HUANG2022100006,WANG2022113021}. This is explained by the chlorophyll reflectance of plants in the red/NIR zones which can be used as an arithmetic combination of bands 4-5 for computing the Normalized Difference Vegetation Index (NDVI) \cite{HUETE2002195} or other vegetation indices \cite{
JIANG20083833,HUETE1988295} as indicators of vegetation health, vigour, and plant phenology, as reported earlier \cite{SOUDANI2012234,Lemenkova202019,1026428,HMIMINA2013145}.

Modelling, reconstruction and prognosis of tropical forest degradation in Africa poses a challenge for both environmental research as well as socio-economic applications, since deforestation affects current land use policy and agro-industrial development of central African countries \cite{TEGEGNE2016312,f13091508}. Besides climate factors, selected human activities, such as timber production, road expansion, logging and manufacturing contribute to deforestation \cite{BRANDT201615,TRITSCH2020106660,Lietal2015}. Another factor includes a high consumption of wood by local population through agriculture and woodcutting \cite{Iloweka}. At the same time, deforestation and unbalanced functioning of forest ecosystems have wider effects and environmental consequences including land degradation, desertification, the loss of wildlife habitat and soil erosion \cite{AQUILAS2022100553}. Analysis of deforestation is possible by comparison of forest areas on older maps with those detected on the remote sensing data \cite{rs12040638,land10101042,Zhangetal2005}, or landscape dynamics evaluated by time series of images \cite{land11091541}. Other approaches include continuous metric to measure the degree of forest degradation \cite{SHAPIRO2021107268}
 
While the techniques of the remote sensing data acquisition and image processing by GIS have demonstrated advances over the recent decades \cite{CHUMA2021100083,Westmanetal1989}, the algorithms of geospatial data modelling are not as straightforward to handle the data automatically. At the same time, the required automation is possible using the programming approaches. A large number of methods focused on forest monitoring by the remotely sensed data have been proposed over the recent years \cite{6350700,773475,PU2019268,6049235,OHMANN20143}. Popular cases include the use of the multi-source data, such as a fusion of LiDAR, radar and optical data, very high-resolution satellite images and unmanned aerial vehicles (UAV) \cite{rs12071087}. The application of land use/land cover LULC change \cite{lambin2003dynamics}, classification and clustering \cite{976841}, and integration of the remote sensing data with reference maps for change detection \cite{rs14164126}. Typically, methods of forest mapping are based on the use of the Landsat imagery \cite{1294510,Zhangetal2006,KIM2014178}, others utilize SAR \cite{rs11242999,858324,7326669,6352093,8035264} or MODIS scenes \cite{HANSEN20082495,ZHANG201974} for land cover classification in forest areas. 

Compared with GIS, the use of programming languages for image processing presents new possibilities in mapping through the extended functionality of algorithms, and has received much attention in geoscience community \cite{TRUJILLOJIMENEZ2022100703,VANSTRIEN2022105462,RODRIGUES2022711,5981114,ROBERTS2022105192}. As previous work shows \cite{LIAQAT2021100587,ASUQUOENOH2022,RAWAT201577}, integrating spatial information from the satellite images using various perceptual tasks such as object recognition, location, classification and interpretation are necessary for disambiguating visually similar land cover classes of the Earth's surface representing the landscapes. Automatic recognition of the diverse land cover classes on various periods facilitates the recognition of the environmental changes. However, misclassifications remain in challenging situations of similar classes with human-based interpretation using supervised classification GIS for image analysis \cite{HUSSAIN2022103117,ABDELSADEK2022815,HEDAYATI202273}. This requires to recognise and classify all possible occurrence of pixels on the scenes automatically, which is a computationally effective task using programming algorithms.

In this paper, we propose an application of R programming language \cite{RCoreTeam} as a means to process and classify satellite images. Specifically, we presented the unsupervised classification of the Landsat-8 OLI/TIRS data and calculation of the Normalized Difference Vegetation Index (NDVI) based on RSToolbox, raster and terra packages. Motivated by recent reports regarding the environmental changes in rainfall tropical forests of central Africa \cite{Wasseige,lescuyer2010chainsaw,Jayathilake,POTAPOV2012106,Shapiro,7326202}, we applied multi-temporal data covering identical selected regions along the middle Congo River Basin, with a time span of 2013 to 2022. The practical aim is to demonstrate the dynamics of the landscapes within a decade using Landsat satellite images. Our methods apply the k-means automated clustering and the calculation of NDVI index using R libraries RStoolbox \cite{RStoolbox}, terra \cite{RTerra}, raster \cite{Hijmans} and auxiliary graphical packages such as RColorBrewer \cite{Neuwirth} and viridis \cite{viridis}. 

%%%%%%%%%%%%%%%%%%%%%%%%%%%%%%%%%%%%%%%%%%
\section{Study Area}

The dataset contains satellite images obtained for the region of Democratic Republic of Congo (DRC) with the three target areas located in the middle region of the Congo River Basin, according to the regional division \cite{MUSHI201964}. The surroundings of the the towns Bumba, Basoko and Kisangani are located on the Congo River, which serves as the main transportation artery in the region (Figure \ref{fig01}). 

Kisangani is the capital of the Tshopo province, situated at the confluence of the rivers Congo, Lualaba, Tshopo, and Lindi. It is strategically located at the crossroads between the eastern and western parts of Congo. Basoko town is situated in the confluence of the Aruwimi tributary into the Congo River. Bumba is a small town and a river port in the Mongala Province, northern part of the DRC. As a part of the Congo River Basin, these cities have a tropical monsoon climate with dry winter period, humid summer-autumn seasons with heavy precipitation and periodically flooded areas \cite{6571272}. Such climate setting, characterised by rainfalls interspersed with dry season climate in Congo Basin create favourable conditions for the distribution of world's largest tropical humid forests \cite{Washington}. 

The importance of the study area is explained by the crucial environmental role of the Congo River Basin. It contains over half of the African tropical moist forests with Salonga National Park being the largest rain forest park in the world \cite{Sayer1992}, and it largely contributes to the Earth's atmosphere circulation on the planetary scale. Besides, the tropical rainforests of Congo are a key factor in producing the oxygen and regulating global temperature through the absorption of solar radiation by forest massifs. Moreover, forest resources offer critical and principal support to the welfare of the rural livelihoods and socioeconomic wellbeing of local population \cite{f13010130} which includes forest-related practices such as forest clearing. As a result, land use and shifting agriculture become the major causes of hydrological flow regime, environmental changes from regional to global scale \cite{4780136}, as well as deforestation along with the increased population density and associated urban growth \cite{Wilkie1994}. 

At the same time, primary forests are essential ecosystems that play a key role in mitigating climate change which raises the need for taking actions on mitigating the effects from deforestation and resilient development in forest degradation \cite{SOMORIN2012288,BELE20151}. Conserving natural resources in Congo tropical rainforests recently launched global initiatives, such as Reducing Emissions from Deforestation and Forest Degradation (\href{https://www.fao.org/redd/overview/en/}{REDD+}) \cite{land11060918,SUFOKANKEU2020102081} or Global Rain Forest Mapping (GRFM) project \cite{1025720}. On a regional and local level, practicing agroforests serve as ecosystem services responding to the climate change mitigation strategy \cite{Batsietal2021}. Nevertheless, Congo River Basin remains the least studied region in Africa and in the global tropics \cite{Washingtonetal2006}. The approach we present here contributes to the development of methods for the environmental analysis of the selected regions of Congo River Basin using remote sensing data and R programming.

\begin{figure}[H]
\centering
	\includegraphics[width=12.0 cm]{Fig_01}
	\caption{Study area with the location of the three target cities in Congo River Basin, DCR. Cartography: Generic Mapping Tools (GMT). Data source: GEBCO/SRTM. Map source: authors.
	\label{fig01}}
\end{figure} 

%%%%%%%%%%%%%%%%%%%%%%%%%%%%%%%%%%%%%%%%%%
\section{Materials and Methods}

This article proposes a general a fully automated method for classifying satellite images using R libraries by machine based \emph{k}-means clustering which requires no prior information and no learning on land cover classes. On one hand, the nature of land cover types is captured by considering the spectral attributes of the pixels in the respective Landsat bands. By carefully analysing these attributes, we reveal and demonstrate that the irregularity of the spectrum in wavelength corresponding to the landscape properties is highly related to the reflectance of of the surface which results in the brightness of pixels on the image. Thus, we applied the computational model of R to effectively capture and visualise this irregularity and transfer it from the frequency optical domain to the geo-information domain.   

\subsection{Research Design}
Within the scope of the current study, the functionality of R is focused on the K-means clustering classification algorithm of RStoolbox to process the satellite images, with aim to visualise and compare the vegetation extent in central Congo River Basin for different years using general workflow approach, Figure \ref{fig:work}. The Landsat satellite images taken on 2013, 2015 and 2022 were used for visualising change detection between land cover or forest types. Thus, the algorithms of R support the operational monitoring of the dynamics of land cover types derived from the images. However, the packages of R can also be applied to additional applications of geospatial data analysis where automated image classification is needed. Examples of such cases include monitoring the desertification using time series of images, analysis of urban growth using retrospective data with overlaid with current maps and actual images, or forest fire monitoring. 

\begin{figure}[H]
\centering
	\includegraphics[width=14.0 cm]{workflow.jpg}
	\caption{Workflow of the study project summarising general steps of the study, including aim and objectives, methodology and results used in a framework of R application for remote sensing data processing, land cover change detection by unsupervised classification (k-means clustering) of Landsat satellite images for selected regions of Congo, D.R.C.
	\label{fig:work}}
\end{figure} 

The functionality of R programming presents a set of modular functions for the automation of image classification applications from Earth observation, spatial and remote sensing data, Figure \ref{fig:Rfunc}. 

\begin{figure}[H]
\centering
	\includegraphics[width=14.0 cm]{R_funct.png}
	\caption{The functionality of R packages used for satellite image processing and classification in this study. Major packages include terra, raster, Rstoolbox, rgdal, rgeos, cluster, graphics, RColorBrewer.
	\label{fig:Rfunc}}
\end{figure} 

The packages of 'terra' and 'raster' includes functionality of pre- and post-processing of both data types (vector and raster), and extensible features for spatial data analysis. Raster data analysis includes reading the information on spectral reflectance of the image bands\hl{, }automated processing of the stacked bands, feature extraction for unsupervised classification using clustering, creating multi-layer \hl{colour composites}, performing map algebra using raster bands, visualising the thematic maps as overlays, classification and basic statistical analysis such as histograms.

R operates with major features of spatial data, such as dimensions, resolution, reading data extent (min/max analysis), plotting the data, manipulating with Coordinate Reference Systems (CRS), reprojecting the data, assigning the CRS and transforming the data. Important function of R 'terra' package consists in creating the SpatRaster objects, which is a class of raster data with a longitude/latitude CRS as spatial parameters and specific geometry of file (cells arranged in rows/columns of the internal grid). Moreover, the 'terra' package enables cartographic visualization of the continuous fields including topographic elevation, vegetation based on the classified land cover types, climate data in spatial context (regional variations of temperature, and precipitation). Besides, R enables to perform data clipping using overlay of spatial data and imagery over their area of interest.

\subsection{Data}

The first step in workflow included data download from the \href{https://glovis.usgs.gov/app}{USGS GloVis} data repository which offers a global coverage of the publicly available satellite images with open access. We selected the three areas of Bumba, Basoko and Kisangani and collected the Landsat-8 OLI/TIRS imagery for the dates 2013 and 2022. The exact Landsat Product IDs are summarised in Table \ref{tab1}. Main characteristics of the imagery are summarised in Table \ref{tab2}. 
\begin{table}[H] 
\caption{Landsat-8 OLI/TIRS Product ID and Scene ID \label{tab1}}
	\newcolumntype{C}{>{\centering\arraybackslash}X}
	\begin{tabularx}{\textwidth}{| p{2.3cm} | p{6.3cm} | p{5.3cm} |}
\toprule
\textbf{Region, year} & \textbf{Landsat Product ID} & \textbf{Landsat Scene ID}\\
\midrule
	Basoko, 2013 & LC08\_L1TP\_177059\_20131216\_20200912\_02\_T1 & LC81770592013350LGN01\\
	Basoko, 2022 & LC08\_L1TP\_177059\_20220208\_20220212\_02\_T1 & LC81770592022039LGN00 \\
	Bumba, 2013 & LC08\_L1TP\_178058\_20131223\_20200912\_02\_T1 & LC81780582013357LGN01 \\
	Bumba, 2022 & LC08\_L1TP\_178058\_20220130\_20220204\_02\_T1 & LC81780582022030LGN00 \\
	Kisangani, 2013 & LC08\_L1TP\_176060\_20130413\_20200913\_02\_T1 & LC81760602013103LGN02 \\
	Kisangani, 2015 & LC08\_L1TP\_176060\_20150113\_20200910\_02\_T1 & LC81760602015013LGN01 \\
	Kisangani, 2022 & LC08\_L1TP\_176060\_20220217\_20220302\_02\_T1 & LC81760602022048LGN00 \\
\bottomrule
\end{tabularx}
\end{table}

\begin{table}[H]
\begin{threeparttable}
\caption{Main technical characteristics of the Landsat-8 OLI/TIRS imagery}
    \centering
   \small
    \setlength\tabcolsep{3pt} % default: 6pt
%    \begin{tabularx}{\linewidth}{| l | l | l | l | l | l | l | l |}
 %   \begin{tabularx}{\linewidth}{| p{1.5in} | p{1.5in} | p{1.5in} | p{1.5in} | p{1.5in} | p{1.5in} | p{1.5in} | p{1.5in} |}
	 \begin{tabular*}{5in}{@{\extracolsep{\fill}}| p{1.5cm} | p{1.4cm} | p{1.4cm} | p{1.4cm} | p{1.4cm} | p{1.6cm} | p{1.6cm} | p{1.6cm} |}
%	\begin{tabularx}{\textwidth}{@{} ll *{6}{Y} @{}}
    \toprule
        \textbf{Parameters} & \textbf{Basoko (2013)} & \textbf{Basoko (2022)} & \textbf{Bumba (2013)} & \textbf{Bumba (2022)} & \textbf{Kisangani (2013)} & \textbf{Kisangani (2015)} & \textbf{Kisangani (2022)} \\ \hline
        \midrule
        Date & 2013-12-16 & 2022-02-08 & 2013-12-23 & 2022-01-30 & 2013-04-13 & 2015-01-13 & 2022-02-17 \\ \hline
        TWRS P. & 177 & 177 & 178 & 178 & 176 & 176 & 176 \\ \hline
        WRS Row & 59 & 59 & 58 & 58 & 60 & 60 & 60 \\ \hline
        WRS Path & 177 & 177 & 178 & 178 & 176 & 176 & 176 \\ \hline
        Cloudiness & 0.00 & 0.00 & 0.04 & 0.00 & 5.85 & 1.36 & 0.48 \\ \hline
        UTM Zone & 34 & 34 & 34 & 34 & 35 & 35 & 35 \\ \hline
        \bottomrule
%    \end{tabularx}\label{tab:01}
\end{tabular*}\label{tab2}
 \begin{tablenotes}
      \small
      \item Abbreviations used in the Table 2: SA -- Sun Azimuth. Date stands for the acquisition date of the scene; GCP Ver. -- GCP Version; TWRS P. -- Target WRS Path; GCP Mod. -- GCP Model; GRMSEM -- Geometric RMSE Model.
    \end{tablenotes}
  \end{threeparttable}
\end{table}

The common technical parameters for all the scenes are as follows: Spacecraft ID -- Landsat 8; Origin: USGS; Station ID -- LGN; Sensor ID -- OLI\_TIRS; Processing level -- L1TP; Collection Number -- 2; Collection category -- T1;  Output format -- GEOTIFF; WRS type -- 2; Datum and Ellipsoid -- WGS84; Data Source Elevation -- GLS2000; GCP Version -- 5.

The used images are categorised by the L1TP Level 1 Precision Terrain (Corrected), which are the DNs in a 16-bit integer format, converted to Top of Atmosphere (TOA) spectral reflectance (Bands 1--9) or radiance (Bands 1--11) using scaling factors \cite{Engebretson}. The L1TP images include radiometric, geometric, and precision correction, and use a Digital Elevation Model (DEM) to correct errors caused by local topographic relief; the accuracy of the precision/terrain-corrected product depends on the Ground Control Points (GCPs) obtained from the GLS2000 dataset, and resolution of the DEM. Thus, the geometric accuracy of the L1TP Landsat images is valiadated using the ground control points (GCP) which are derived from the GLS2000 library. These data are used for computing the RMSE of the geometric residuals on the terrain-corrected images, \cite{Engebretson}. 

The additional image for 2015 was selected for Kisangani due to the higher cloud coverage in 2013. The geographic map in Figure \ref{fig01} has been plotted using the Generic Mapping Tools (GMT) defined by scripts following the existing methodology \cite{Lemenkova202212,Lemenkova202217,Lemenkova2022b}. The study area encompasses the middle Congo river basin with the three specific towns located in consecutive order southwards: Bumba, Basoko, Kisangani. The collected data were inspected for metadata using R Listing \ref{lst:01}. 

\begin{lstlisting}[language=r, caption=R code used for data inspection by rGDAL and raster libraries; here a case of Basoko town applied for all other images, basicstyle=\scriptsize, label={lst:01}]
library(rgdal)
library(raster)
# Set up working directory
setwd("/Users/polinalemenkova/Documents/R/53_UC_k_means/Basoko_LC08_L1TP_177059_20131216_20200912_02_T1")
# Importing data: Landsat OLI/TIRS image for Basoko, Congo (2013):
Landsat_Basoko2013 <- list.files("/Users/polinalemenkova/Documents/R/53_UC_k_means/Basoko_LC08_L1TP_177059_20131216_20200912_02_T1")
# Printing the list
list.files()
# Reading in files as SpatialGridDataFrame objects :
BasokoBand1 <- readGDAL("LC08_L1TP_177059_20131216_20200912_02_T1_B1.TIF")
BasokoBand2 <- readGDAL("LC08_L1TP_177059_20131216_20200912_02_T1_B2.TIF")
BasokoBand3 <- readGDAL("LC08_L1TP_177059_20131216_20200912_02_T1_B3.TIF")
BasokoBand4 <- readGDAL("LC08_L1TP_177059_20131216_20200912_02_T1_B4.TIF")
BasokoBand5 <- readGDAL("LC08_L1TP_177059_20131216_20200912_02_T1_B5.TIF")
BasokoBand6 <- readGDAL("LC08_L1TP_177059_20131216_20200912_02_T1_B6.TIF")
BasokoBand7 <- readGDAL("LC08_L1TP_177059_20131216_20200912_02_T1_B7.TIF")
# Visualizing and checking the class of random bands
plot(BasokoBand1)
res(BasokoBand1)
class(BasokoBand5)
## [1] "SpatialGridDataFrame"
## attr(,"package")
## [1] "sp"
slotNames(BasokoBand5)
# [1] "data"        "grid"        "bbox"        "proj4string"
summary(BasokoBand5@data)
# Min.   : 6637
# 1st Qu.:14955
# Median :15669
# Mean   :15752
# 3rd Qu.:16579
# Max.   :27433
# NA's   :17221239
# Check the topology of the class
str(BasokoBand5@grid)
#Formal class 'GridTopology' [package "sp"] with 3 slots
# ..@ cellcentre.offset: Named num [1:2] 696300 43500
#  .. ..- attr(*, "names")= chr [1:2] "x" "y"
#  ..@ cellsize         : num [1:2] 30 30
#  ..@ cells.dim        : int [1:2] 7591 7741
# Check up spatial dimension
BasokoBand5@grid@cellcentre.offset
# x      y
# 696300  43500
BasokoBand5@grid@cellsize
# [1] 30 30
BasokoBand5@bbox
#      min    max
# x 696285 924015
# y 43485 275715
BasokoBand5@proj4string
# generating RasterLayers of Landsat bands
BasokoB1 <- raster("LC08_L1TP_177059_20131216_20200912_02_T1_B1.TIF")
BasokoB2 <- raster("LC08_L1TP_177059_20131216_20200912_02_T1_B2.TIF")
BasokoB3 <- raster("LC08_L1TP_177059_20131216_20200912_02_T1_B3.TIF")
BasokoB4 <- raster("LC08_L1TP_177059_20131216_20200912_02_T1_B4.TIF")
BasokoB5 <- raster("LC08_L1TP_177059_20131216_20200912_02_T1_B5.TIF")
BasokoB6 <- raster("LC08_L1TP_177059_20131216_20200912_02_T1_B6.TIF")
BasokoB7 <- raster("LC08_L1TP_177059_20131216_20200912_02_T1_B7.TIF")
# Creating a facetted stack from the Landsat OLI-TIRS bands:
image <- stack(BasokoB1, BasokoB2, BasokoB3, BasokoB4, BasokoB5, BasokoB6, BasokoB7)
# Visualizing the image with separated bands:
plot(image)
# Check metadata:
nlayers(image)
res(image)
# Creating the object of RasterLayer class for one random band (here: Band 2)
Basoko2013_B2 <- raster(Landsat_Basoko2013[2])
Basoko2013_B2
# Plotting one randon band
plot(Basoko2013_B2,
     main = "Landsat OLI/TIRS 2013 band 2\nBasoko, central Congo",
     col = gray(0:100 / 100))
\end{lstlisting}

We explored the 'rgdal' library of R to retrieve the coordinates, check the topology and spatial dimension of the dataset stored in the SpatialGridDataFrame class. A topology refers to the arrangement of the images that defines the coincident geometry of the feature classes from the Landsat dataset. It reads the contextual information from the structure data based on spatial relationships: the adjacency of the coordinates and connectivity of overlapping images. The RasterLayers were generated as the arrays of pixels from the individual bands of the Landsat scenes. Afterwards, the facetted stack was created using the Landsat OLI-TIRS bands by stack() function. To this end, the 'raster' R library was applied for concatenating the multiple bands into a single bitmap image which uses the data on the origin of the coordinates of each band. The resulting image was visualized by the plot() function of R. The number of layers and spatial dimension of in the new multi-layer object of R were inspected by 'raster' library. The resolution of all the layers was 30 m except for the panchromatic band (15 m) and the two last TIRS layers with a coarser resolution 100 m. The selected band was plotted in a grayscale monochrome visualization. The data analysis and inspection has been performed using Listing \ref{lst:01}.

\subsection{Plotting band composites}

We explored the structure of Landsat bands by inspecting the metadata. The image of all the bands is visualized as a representation of the monochrome multi-facetted plot in Figure \ref{fig02}. The most of the multispectral bands have 30-m resolution, while the panchromatic band (B8) has a higher resolution (15 m) and the two TIRS (B10 and B11) have a coarse 100-m resolution. Generating colour composites by combination of various band triplets provides more information which can be effectively exploited using libraries 'raster' and 'terra' of R. Here the idea of the machine-assisted classifier by R packages consists in the identification of spectral reflectance of pixels constituting the 11 bands for thematic mapping. 

The sensor instruments of Landsat-8 OLI-TIRS record more channels of data than are actually needed for vegetation applications. For instance, for NDVI calculation we only need the Red and NIR bands (4 and 5, respectively). However, combining the bands in various RGB triplets enables to highlight diverse objects on the surface, to discriminate and capture information on a wider variety of objects. Thus, the composition of red, green and blue colours creates the images with varied hue, saturation and intensity representation owing to the pixel brightness as spectral reflectance of the objects and surfaces on the Earth. Therefore, the bands were mixed to form various colour composites and inspect the land cover types. The selected and most representative band composites for the three main studied regions were plotted as natural and false colour composites, respectively, which were generated using library 'raster'. The relevant code snippet is given for the town of Basoko, and was applied similarly for the rest of the Landsat-8 OLI-TIRS scenes with changed parameters, Listing \ref{lst:02}. 

\begin{figure}[H]
\makebox[\textwidth][r]{\includegraphics[width=1.1\textwidth]{Fig_02.jpg}}%
	\caption{Composite of monochrome 11 bands: Landsat-8 OLI/TIRS C1 for Basoko region (Arowumi and Congo rivers). Multi-spectral image 'LC81770592022039LGN00': Ultra Blue (B1), Blue (B2), Green (B3), Red (B4), Near Infrared (NIR) (B5), Shortwave Infrared (SWIR) 1 (B6), Shortwave Infrared (SWIR) 2 (B7), Panchromatic (B8), Cirrus (B9), Thermal Infrared (TIRS) 1 (B10), TIRS2 (B11). Plotting: R.}
	\label{fig02}
\end{figure} 

\begin{lstlisting}[language=r, caption=R code used for colour composites by libraries raster and
terra; here a case of Bumba town applied likewise for all the other images, label={lst:02},basicstyle=\scriptsize]
setwd("/Users/polinalemenkova/Documents/R/53_UC_k_means/Bumba_LC08_L1TP_178058_20131223_20200912_02_T1")
# Importing data
Landsat_Bumba2013 <- list.files("/Users/polinalemenkova/Documents/R/53_UC_k_means/Bumba_LC08_L1TP_178058_20131223_20200912_02_T1")
# Printing the list
list.files()
# creating SpatRaster object
landsat <- rast(Landsat_Bumba2013)
# check properties
landsat
# combining bands for natural colour composite, here using bands 4, 3 and 2.
landsatRGB <- landsat[[c(4,3,2)]]
# check up metadata
landsatRGB
# visualising  natural colour composite on screen
plotRGB(landsatRGB, r=1, g=2, b=3, axes=FALSE, stretch="lin")
# plotting and saving the file
plot(landsatRGB, col=colors, font.main = 1, main = "NDVI for Landsat-8 OLI/TIRS C1 image LC08_L1TP_178058_20131223_20200912_02_T1: Bumba, Congo", cex.main=0.9, axes=FALSE)
# combining bands for false colour composite, here using bands 5, 4 and 3.
landsatFCC <- landsat[[c(5,4,3)]]
# visualising false colour composite on screen
plotRGB(landsatFCC, r=1, g=2, b=3, axes=FALSE, stretch="lin")
# plotting and saving the file
plot(landsatRGB, col=colors, font.main = 1, main = "NDVI for Landsat-8 OLI/TIRS C1 image LC08_L1TP_177059_20131216_20200912_02_T1_B: Basoko, central Congo", cex.main=0.9, axes=FALSE)
\end{lstlisting}

The multi-facetted plot of the individual bands of the Landsat OLI/TIRS image has been plotted using libraries 'raster' and 'terra' (version 1.2.11), by the code in Listing \ref{lst:03}. Here the case is given for the image covering Basoko (2022), applied for other scenes. 

\begin{lstlisting}[language=r, caption=R code used for the monochrome individual bands of the image Landsat OLI/TIRS, basicstyle=\scriptsize, label={lst:03}]
library(raster)
library(terra)
# Setup working directory
setwd("/Users/polinalemenkova/Documents/R/53_UC_k_means/Basoko_LC08_L1TP_177059_20220208_20220212_02_T1")
b1 <- rast('LC08_L1TP_177059_20220208_20220212_02_T1_B1.TIF') # Ultra Blue/Aerosol
b2 <- rast('LC08_L1TP_177059_20220208_20220212_02_T1_B2.TIF') # Blue
b3 <- rast('LC08_L1TP_177059_20220208_20220212_02_T1_B3.TIF') # Green
b4 <- rast('LC08_L1TP_177059_20220208_20220212_02_T1_B4.TIF') # Red
b5 <- rast('LC08_L1TP_177059_20220208_20220212_02_T1_B5.TIF') # Near Infrared NIR
b6 <- rast('LC08_L1TP_177059_20220208_20220212_02_T1_B6.TIF') # Green
b7 <- rast('LC08_L1TP_177059_20220208_20220212_02_T1_B7.TIF') # Red
b8 <- rast('LC08_L1TP_177059_20220208_20220212_02_T1_B8.TIF') # Panchromatic
b9 <- rast('LC08_L1TP_177059_20220208_20220212_02_T1_B9.TIF') # Cirrus
b10 <- rast('LC08_L1TP_177059_20220208_20220212_02_T1_B10.TIF') # TIRS 1
b11 <- rast('LC08_L1TP_177059_20220208_20220212_02_T1_B11.TIF') # TIRS 2
# Defining the color palette
colors <- kovesi.cyclic_grey_15_85_c0(100)
# Plotting 11 individual layers (raw bands) of the Landsat-8 image.
par(mfrow = c(4, 3), mar=c(0, 5, 0, 0), mai = c(0.1, 0.1, 0.1, 0.1)) # c(bottom, left, top, right)
plot(b1, main = "Ultra Blue (B1)", col = colors, axes = FALSE, legend = TRUE)
plot(b2, main = "Blue (B2)", col = colors, axes = FALSE, legend = TRUE)
plot(b3, main = "Green (B3)", col = colors, axes = FALSE, legend = TRUE)
plot(b4, main = "Red (B4)", col = colors, axes = FALSE, legend = TRUE)
plot(b5, main = "NIR (B5)", col = colors, axes = FALSE, legend = TRUE)
plot(b6, main = "SWIR (B6)", col = colors, axes = FALSE, legend = TRUE)
plot(b7, main = "SWIR (B7)", col = colors, axes = FALSE, legend = TRUE)
plot(b8, main = "Panchromatic (B8)", col = colors, axes = FALSE, legend = TRUE)
plot(b9, main = "Cirrus (B9)", col = colors, axes = FALSE, legend = TRUE)
plot(b10, main = "TIRS 1 (B10)", col = colors, axes = FALSE, legend = TRUE)
plot(b11, main = "TIRS 2 (B11)", col = colors, axes = FALSE, legend = TRUE)
\end{lstlisting}

The generated images reveal spatial details which can be better discerned in the intensity of pixels in false colour composites for water areas and more sharpened representation of urban areas in the natural colour composites showing spatial details in the three target areas: Bumba, Basoko and Kisangani, Figure \ref{fig03}. The colour products for all the three regions were prepared using the combination of bands 4-3-2 (natural) and 5-4-3 (false) with the spatial resolution of the visual bands selected as three triplets of the original multispectral bands of the Landsat-8 OLI/TIRS. The new mixture from the RGB bands is converted to the colours with the intensity channel corresponding to the resolution of the original bands (30 m). A hue, saturation and intensity transformation of each of the triplets are processed and resampled by R package 'terra' to form the output image, Figure \ref{fig03}.

The demonstrated above scripts show the key advantage of R which consists in the optimization of the workflow of remote sensing data processing. Moreover, it illustrates its advanced analysis capabilities aimed at processing raster data by scripting. Scripts can be re-used with modified lines of code for the next images as a consecutive work, rather than repeated as a whole process of image analysis for each image, as in GIS, which is a fundamentally different approach compared to GIS. This is especially actual for processing big Earth observation data which play a key role in geoscience. Rapid development of remote sensing technologies resulted in the increased amount of available satellite images, which requires advanced technologies for processing these large massifs of spatial data. In contrast to the traditional software, programming languages provide an effective solution to handle such volumes of spatial data by scripting. R, as an example of the high-level programming language, effectively solves the problem of optimization of image processing.

\begin{figure}[H]
\makebox[0.90\linewidth][r]{%
	\begin{subfigure}[b]{.6\textwidth}
		\centering
			\includegraphics[width=0.87\textwidth]{Fig_03a.jpg}
		\caption{Bumba: natural colour composite, 2013}
	\end{subfigure}%
	\begin{subfigure}[b]{.6\textwidth}
		\centering
		\includegraphics[width=0.87\textwidth]{Fig_03b.jpg}
		\caption{Bumba: false colour composite: 2013.}
	\end{subfigure}%
}\\
\vfill \vspace{1mm}
\makebox[0.90\linewidth][r]{%
	\begin{subfigure}[b]{.6\textwidth}
		\centering
			\includegraphics[width=0.87\textwidth]{Fig_03c.jpg}
		\caption{Basoko: natural colour composite: 2013.}
	\end{subfigure}%
	\begin{subfigure}[b]{.6\textwidth}
		\centering
		\includegraphics[width=0.87\textwidth]{Fig_03d.jpg}
		\caption{Basoko: false colour composite: 2013}
	\end{subfigure}%
}\\
\vfill \vspace{1mm}
\makebox[0.90\linewidth][r]{%
	\begin{subfigure}[b]{.6\textwidth}
		\centering
			\includegraphics[width=0.87\textwidth]{Fig_03e.jpg}
		\caption{Kisangani: Natural colour composite: 2013.}
	\end{subfigure}%
	\begin{subfigure}[b]{.6\textwidth}
		\centering
		\includegraphics[width=0.87\textwidth]{Fig_03f.jpg}
		\caption{Kisangani: False colour composite: 2013}
	\end{subfigure}%

}
\caption{Natural (bands 4-3-2) and false colour (bands 5-4-3) composites of the Landsat 8 OLI/TIRS image for the three target areas: Bumba, Basoko and Kisangani. Mapping: RStudio}\label{fig03}
\end{figure}

\subsection{k-means Clustering}
The \emph{k}-means clustering of R package RasterStack is a method by which pixels are assigned to spectral classes by machine-based partition without prior knowledge of those classes. It is a straightforward way to produce a map of land cover classes using automatic approach of image classification. The algorithm has been performed using image-partitioning algorithm that infers the labels of the land cover classes by analysing intra-class similarity and, vice versa, the contrast between the classes. This method presents a machine learning approach for assigning pixels into the defined groups (or coherent clusters), according to their intensity and brightness, Listing \ref{lst:04}. Clustering and classification of data is a commonly accepted practice for land cover mapping aimed at modelling changes in a forest structure by examining land-use categories  \cite{Manakos,WILKIE1988307,4780136,rs11242999,land11060787}.

\begin{lstlisting}[language=r, caption=R code used for running the unsupervised classification by k-means clustering; here a case of Basoko town applied for all other images, basicstyle=\scriptsize, label={lst:04}]
# Set up working directory
setwd("/Users/polinalemenkova/Documents/R/53_UC_k_means/Basoko_LC08_L1TP_177059_20220208_20220212_02_T1")
# Importing data
Landsat_Basoko2022 <- list.files("/Users/polinalemenkova/Documents/R/53_UC_k_means/Basoko_LC08_L1TP_177059_20220208_20220212_02_T1")
# Printing the list
list.files()
# Creating the object of RasterLayer class for one band
Basoko2022_B2 <- raster(Landsat_Basoko2022[2])
Basoko2022_B2
# Generating the color palette table
colors <- kovesi.linear_grey_10_95_c0(100)
# Plotting one band
plot(Basoko2022_B2,
    main = "Landsat OLI/TIRS 2022 band 2 Basoko, Congo \nLC08_L1TP_177059_20220208_20220212_02_T1", font.main=2, cex.main = 0.90, axes = FALSE, box = FALSE, legend = FALSE, col = colors)
# Stacking the data to create a RasterStack. 
Landsat_Basoko2022_stack <- stack(Landsat_Basoko2022)
# Turning a stack into a brick. 
# Landsat_Basoko2022_brick <- brick(Landsat_Basoko2022_stack)
Landsat_Basoko2022_brick <- brick(Basoko2022_B2)
# Viewing brick attributes
Landsat_Basoko2022_brick
# Plotting the RGB triplet from the RasterBrick
olpar <- par(no.readonly = TRUE) # back-up par
par(mfrow=c(1,1))
plotRGB(Landsat_Basoko2022_brick)
## Running the classification
set.seed(25)
unC <- unsuperClass(Landsat_Basoko2022_brick, nSamples = 100, nClasses = 10, nStarts = 5)
unC
# Creating the color palette table
colors <- kovesi.diverging_rainbow_bgymr_45_85_c67(10)
# Plotting a map
plot(unC$map, main = "K-means Clustering for Landsat-\hl{8} OLI/TIRS C1 image of Basoko, Congo  \nL1TP_177059_20220208_20220212_02_T1 (2022)", font.main=2, cex.main = 0.95, col = colors, axes = FALSE, box = FALSE, legend = FALSE)
# Adding legend
legend("topleft", legend=c("1", "2", "3", "4", "5", "6", "7", "8", "9", "10"), fill = colors, title = "Classes", horiz = FALSE,  bty = "n", text.font=3, ncol=1)
\end{lstlisting}

The image for classification was generated using stacking the data which enabled to create the RasterStack as a collection of RasterLayer objects. In turn, they were generated from the raster layers of the selected bands with the same spatial extent and resolution, that is, 30 m for the most of the Landsat-8 OLI/TIRS bands. Afterwards, RasterStack was merged into the RasterBrick, which is a multi-layer raster object created from the multi-layer files of RasterStack, made up of Landsat bands. The similarity between the RasterStack and RasterBrick classes consists in the fact that both are generated with a stack of the raster Landsat bands. However, the processing time for the RasterBrick is shorter in R environment, compared to the RasterStack, because RasterBrick is less flexible and is merged already as a single file. The procedure and the output of the \emph{k}-means clustering is presented in Figure \ref{fig04}, which shows the console of the RStudio with statistical information on cluster centroids and sum of squares for each processed map, respectively. The cluster map on the right of the print screen depicts the class memberships of the pixels, in this case we give the example of the Basoko town on the year 2013, Figure \ref{fig04}. 

\begin{figure}[H]
\centering
	\includegraphics[width=15.0 cm]{Fig_04.jpg}
	\caption{Executing the \emph{k}-means clustering classification algorithm in RStudio using the RStoolbox package. Parameters of cluster centroids and sum of squares by cluster are displayed in the console of R. Spatial parameters of the output map are listed. Here the example of Landsat image of the Basoko.
	\label{fig04}}
\end{figure} 

Thus, the automatic clustering analysis by R has been performed without prior knowing of class membership of those pixels using a principle of\hl{ }partitioning\hl{ of the }raster file based on the generalization of the matrix of cells. Based on the automatic analysis of pixel brightness, the pixels of a given symbol were partitioned into the same spectral class which belongs to one of the 10 clusters. The advantage of this approach is that it works fast and straightforward in operating complex scenes. The clusters were identified by associating the pixels from the images with the available information of spectral reflectance and presented as the output maps in Figure \ref{fig05} (Bumba), Figure \ref{fig06} (Basoko) and Figure \ref{fig07} (Kisangani). 

The unsupervised learning algorithm of R incorporates the parameter estimation for pixels under partitioning framework using cell's brightness that corresponds to the spectral reflectance of the objects on a matrix of the raster scene. The algorithm of \emph{k}-means clustering is computationally demanding by R, yet it is crucial to the concept of remote sensing due to the effectiveness and importance of the output in satellite image processing. Thus, land cover types are not known beforehand and defined by the machine method of image analysis, based on the spectral signatures of the objects partitioned into the series of clusters. The data were divided automatically into the ten clusters using the search for spectral class for each pixel in the image. 
Using this approach, we defined 10 classes in a process of the unsupervised classification approach to produce the thematic maps with discriminated vegetation and land cover classes using the existing classification adopted from previous studies \cite{Bwangoy}: swamp forest; humid tropical forest; secondary forest; urban areas and impervious surfaces; inundated grassland; deciduous forest; shrubland and grassland; cropland and mosaic forest; savannah with sparse trees; water bodies. 

\subsection{NDVI calculation}
The NDVI is a key vegetation indicators to describe vegetation dynamics and vigour. It is effectively applied to environmental management and monitoring the responses of tropical rainforests to climate change and anthropogenic activities. Therefore, the use of the NDVI is crucial to reveal the health of the terrestrial ecosystems in central Africa which are threatened by the deforestation. As a robust traditional estimation method to determine vegetation coverage, NDVI is deservedly applied in a variety of the remote sensing studies due to its effectiveness and reliability. Besides, the visualization of the NDVI is indispensable in areas where fieldwork and ground points collections is hardly accessible, such as tropical African regions of central Congo. Therefore, we adapted the algorithm of NDVI for extracting the information on vegetation for knowledge on plant coverage in the three target areas in Congo Basin. The NDVI has been generated using composite band arithmetic operations by ratio of difference and the sum of the reflected NIR and visible Red wavelength. The NDVI approach is grounded on the calculation of ratios in the reflectance between the Red and NIR bands of the image (Band 4 and 5), as arithmetical variables using the existing formula $NDVI=(NIR-Red)/(NIR+Red)$. The NDVI computation has been performed using multi-band manipulation by R library terra, Listing \ref{lst:05}. 

\begin{lstlisting}[language=r, caption=R code using library terra used for calculation the NDVI; here a case of Basoko town (2013) with the same principle applied for all other images, basicstyle=\scriptsize, label={lst:05}]
# Set up working directory
setwd("/Users/polinalemenkova/Documents/R/53_UC_k_means/Basoko_LC08_L1TP_177059_20131216_20200912_02_T1")
# Importing data
filenames <- paste0('LC08_L1TP_177059_20131216_20200912_02_T1_B', 1:7, ".tif")
filenames
landsat <- rast(filenames)
landsat
vi <- function(img, k, i) {
  bk <- img[[k]]
  bi <- img[[i]]
  vi <- (bk - bi) / (bk + bi)
  return(vi)
}
# plotting the NDVI. For Landsat NIR = 5, red = 4.
ndvi <- vi(landsat, 5, 4)
colors <- brewer.pal(10, "RdYlGn")
plot(ndvi, col=colors, font.main = 1, main = "NDVI for Landsat-8 OLI/TIRS C1 image LC08_L1TP_177059_20131216_20200912_02_T1: Basoko, Congo DC (2013)", cex.main=0.95, axes = FALSE, legend = FALSE)
legend(, x = "topleft", inset = 0.13, legend=c("0.1", "0.2", "0.3", "0.4", "0.5", "0.6", "0.7", "0.8", "0.9", "1.0"), fill = colors, title = "Classes", horiz = FALSE,  bty = "n", text.font=3, ncol=1)
\end{lstlisting}

Calculation of the NDVI enabled to exploit the information derived from spectral reflectance regarding vegetation coverage. This is enabled due to the parameters of the ratios of NIR/Red spectral bands which lessen the influence from the ground relief, rocks and soils. Thus, the methodological advantage of the NDVI consists in the reinforcement of the slight differences and nuances in spectral reflectance characteristics of vegetation. This enables to better distinct forest covered areas, discriminate various vegetation classes and discern them from all other land cover types. Using pairwise representations of the 9-year time series of Landsat OLI/TIRS (2013 to 2022), we calculated the NDVI for the three target regions along the Congo river: Bumba, Basoko and Kisangani. Here we computed the values of NDVI for each couple of images, respectively, to show the dynamics of the vegetation growth in middle Congo, as shown in Figure \ref{fig08}. The NDVI is computed for each pixel, with brighter green hues indicating more dense and healthy plant canopy, and contrariwise. 

Apparently, the combination of a higher reflectance in NIR and a lower reflectance in Red bands results in a higher NDVI values in general, that is, around 1, which represents spectral signature of vegetation. Such values refer to the healthy plant coverage, typical for the tropical rainforest and other regions of dense canopy, while bank areas of Congo river, bare land, sparsely vegetated areas and urban spaces near Bumba, Basoko and Kisangani have lower NDVI values, respectively. For instance, dark brown to black areas of the NDVI with values around zero are non-vegetated areas and water areas have negative values below \hl{zero}. Middle values indicate various types of shrubland, riparian plant communities along the river Congo and tributaries. Thus, the general approach is that in the NDVI the scale is ranging from -1 to 1 with positive values from 0 to 1 indicating the terrain areas, while water surfaces have values below 0. Higher NDVI values (ca. from 0.7 to 1.0) are indicated by the healthy vegetation, while lower values (from 0.0 to 0.3) represent the non-vegetated areas, correspondingly. 

%%%%%%%%%%%%%%%%%%%%%%%%%%%%%%%%%%%%%%%%%%
\section{Results and Discussion}

\subsection{Image classification}

Figure \ref{fig07} shows changes in land cover types in Bumba region for 9-year time according to the k-means classification by R.

\begin{figure}[H]
\makebox[0.90\linewidth][r]{%
	\begin{subfigure}[b]{.6\textwidth}
		\centering
			\includegraphics[width=0.85\textwidth]{Fig_05a.jpg}
		\caption{Monochrome image for Bumba region: 2013.}
	\end{subfigure}%
	\begin{subfigure}[b]{.6\textwidth}
		\centering
		\includegraphics[width=0.95\textwidth]{Fig_05b.jpg}
		\caption{Classified scene using \emph{k}-means clustering: 2013.}
	\end{subfigure}%
}\\
\vfill \vspace{1mm}
\makebox[0.90\linewidth][r]{%
	\begin{subfigure}[b]{.6\textwidth}
		\centering
			\includegraphics[width=0.85\textwidth]{Fig_05c.jpg}
		\caption{Monochrome image for Bumba region: 2022.}
	\end{subfigure}
	\begin{subfigure}[b]{.6\textwidth}
		\centering
		\includegraphics[width=0.95\textwidth]{Fig_05d.jpg}
		\caption{Classified scene using \emph{k}-means clustering: 2022.}
	\end{subfigure}%

}
\caption{Bumba town and surroundings of the Congo River: raw Landsat-8 OLI/TIRS scene for year 2013 and 2022 and the classified outputs by \emph{k}-means clustering. Mapping: RStudio. Source: authors.}\label{fig05}
\end{figure}

Pairwise examination of the classified images allows to examine the changes in forest structure which are caused by the anthropogenic activities, including urban growth, and environmental drivers initiated by climate change. To this end, we present the results of the classification of the three target regions showing changes in tropical moist forest along the flow of Congo River, as shown in Figure \ref{fig05}, Figure \ref{fig06} and Figure \ref{fig07} for Bumba, Basoko and Kisangani, respectively. 

\begin{figure}[H]
\makebox[0.90\linewidth][r]{%
	\begin{subfigure}[b]{.6\textwidth}
		\centering
			\includegraphics[width=0.85\textwidth]{Fig_06a.jpg}
		\caption{Monochrome image for Basoko region: 2013.}
	\end{subfigure}%
	\begin{subfigure}[b]{.6\textwidth}
		\centering
		\includegraphics[width=0.95\textwidth]{Fig_06b.jpg}
		\caption{Classified scene using \emph{k}-means clustering: 2013.}
	\end{subfigure}%
}\\
\vfill \vspace{1mm}
\makebox[0.90\linewidth][r]{%
	\begin{subfigure}[b]{.6\textwidth}
		\centering
			\includegraphics[width=0.85\textwidth]{Fig_06c.jpg}
		\caption{Monochrome image for Basoko region: 2022.}
	\end{subfigure}
	\begin{subfigure}[b]{.6\textwidth}
		\centering
		\includegraphics[width=0.95\textwidth]{Fig_06d.jpg}
		\caption{Classified scene using \emph{k}-means clustering: 2022.}
	\end{subfigure}%
}
\caption{Basoko town (Tshopo Province) and surroundings in the confluence of the Aruwimi River tributary into the main stream of the Congo River: raw Landsat-8 OLI/TIRS scene for year 2013 and 2022 and the classified outputs by \emph{k}-means clustering. Mapping: RStudio. Source: authors.}\label{fig06}
\end{figure}

Classes of the land cover types were defined as follows: Swamp forest; Humid tropical forest; Secondary forest; Urban areas and impervious surfaces; Inundated grassland; Deciduous forest; Shrubland and grassland; Cropland and mosaic forest; Savannah with sparse trees; Water bodies. The model was done by comparing the results of the vegetation types classified based on the Landsat imagery in 2013 and 2022. The additional image was selected for the Kisingani for 2015 due to the cloudiness of the image in 2013. As demonstrated on the resulting plots, the pixel-based models use 10 land-cover categories. The integration of the two Landsat OLI/TIRS images for land cover mapping by R language is a challenging and promising approach which can be recommended for future similar studies on rainforest mapping. The final outputs as classified maps contained 10 land cover classes recorded from the 30 m raster grids equivalent to the resolution of the initial input Landsat OLI/TIRS data. 

Land cover classification using R approach for processing remote sensing data by RStoolbox is a promising and effective approach. However, a pixel-based data processing of the satellite images results in the grid structure of the landscape with resolution corresponding to the initial data. For more details data, Sentinel scenes can be considered as well for a more fine-resolution classification. The continued increase of the volume of spatial Earth observation data from the satellite images and novel directions in image analysis techniques using programming tools encourage the development of object based approaches. Therefore, technical variations in the classification process can be omitted by applying the object based image analysis for detecting the landscape structure. This can be considered as a further development of the R programming methods.

\begin{figure}[H]
\makebox[0.90\linewidth][r]{%
	\begin{subfigure}[b]{.6\textwidth}
		\centering
			\includegraphics[width=0.88\textwidth]{Fig_07a.jpg}
		\caption{Monochrome image for Kisangani area, B5 (NIR): 2013.}
	\end{subfigure}%
	\begin{subfigure}[b]{.6\textwidth}
		\centering
		\includegraphics[width=0.95\textwidth]{Fig_07b.jpg}
		\caption{Classified scene using \emph{k}-means clustering: 2013.}
	\end{subfigure}%
}\\
\vfill \vspace{1mm}
\makebox[0.90\linewidth][r]{%
	\begin{subfigure}[b]{.6\textwidth}
		\centering
			\includegraphics[width=0.97\textwidth]{Fig_07c.jpg}
		\caption{Classified scene using \emph{k}-means clustering: 2015.}
	\end{subfigure}%
	\begin{subfigure}[b]{.6\textwidth}
		\centering
		\includegraphics[width=0.96\textwidth]{Fig_07d.jpg}
		\caption{Classified scene using \emph{k}-means clustering: 2022.}
	\end{subfigure}%
}
\caption{Kisangani city (Tshopo Province) and region in the junction of the Tshopo and Lindi tributaries into the Congo River: raw Landsat-8 OLI/TIRS scene for year 2013, 2015 and 2022 and the classified outputs by \emph{k}-means clustering. Mapping: RStudio. Source: authors.}\label{fig07}
\end{figure}

\subsection{NDVI}

Three examples of the NDVI computed from the Red/NIR bands of the Landsat OLI-TIRS images have heterogeneous landscapes and vegetation coverage, Figure \ref{fig08}. 

\begin{figure}[H]
\makebox[0.90\linewidth][r]{%
	\begin{subfigure}[b]{.6\textwidth}
		\centering
			\includegraphics[width=0.87\textwidth]{Fig_08a.jpg}
		\caption{Bumba region: 2013.}
	\end{subfigure}%
	\begin{subfigure}[b]{.6\textwidth}
		\centering
		\includegraphics[width=0.87\textwidth]{Fig_08b.jpg}
		\caption{Bumba region: 2022}
	\end{subfigure}%
}\\
\vfill \vspace{1mm}
\makebox[0.90\linewidth][r]{%
	\begin{subfigure}[b]{.6\textwidth}
		\centering
			\includegraphics[width=0.87\textwidth]{Fig_08c.jpg}
		\caption{Basoko region: 2013.}
	\end{subfigure}%
	\begin{subfigure}[b]{.6\textwidth}
		\centering
		\includegraphics[width=0.87\textwidth]{Fig_08d.jpg}
		\caption{Basoko region: 2022.}
	\end{subfigure}%
}\\
\vfill \vspace{1mm}
\makebox[0.90\linewidth][r]{%
	\begin{subfigure}[b]{.6\textwidth}
		\centering
			\includegraphics[width=0.87\textwidth]{Fig_08e.jpg}
		\caption{Kisangani region: 2013}
	\end{subfigure}%
	\begin{subfigure}[b]{.6\textwidth}
		\centering
		\includegraphics[width=0.87\textwidth]{Fig_08f.jpg}
		\caption{Kisangani region: 2022}
	\end{subfigure}%

}
\caption{NDVI computed from 5 and 4 bands of the Landsat-8 OLI/TIRS scenes for year 2013 and 2022 for the three key areas: Bumba, Basoko, Kisangani. Mapping: RStudio. Source: authors.}\label{fig08}
\end{figure}

In the NDVI maps, healthy vegetated areas are bright, water is light to dark brown, and soils and bare ground areas are beige and bluish to slate grey\hl{, Figure} \ref{fig08}. Denoting those shades and hues are understood from the inspection of the relevant land cover types in previous maps. The NDVI is used for the assessment of various vegetation along the Congo flow where the environmental effects from the hydrology on plant growth are more distinct. Moreover, the significance of the NDVI for the environmental monitoring is that it well corresponds to the crop biomass accumulation, as well as to the botanical parameters, for example, such as leaf chlorophyll levels, LAI, and active radiation absorbed by a crop canopy through a photosynthesis process. 

Figure \ref{fig08} demonstrates the three examples of the NDVI maps derived from the three dates of Landsat imagery in Bumba, Basoko and Kisangani regions, respectively. Here the values -1.0 to 0.1 signify water areas (Congo river and tributaries); values 0.1 to 0.2 represent river banks; values 0.2 to 0.3 belong to the class of the non-vegetated areas; pixels with values 0.3 to 0.4 represent seasonally flooded and inundated areas, wetlands and swamps; 0.4 to 0.5 are savannah with sparse trees; pixels of class 0.5 to 0.6 are primarily irrigated agricultural fields with the diverse types of crops; bright pixels of class 0.6 to 0.7 are secondary forests; values of 0.7 to 0.8 are densely vegetated lands and mosaic forests; finally, the areas with NDVI values greater than 0.8 are dense humid tropical forests. 

%%%%%%%%%%%%%%%%%%%%%%%%%%%%%%%%%%%%%%%%%%
\section{Conclusions}
%We have presented 
In this study, we demonstrated that the methods of geospatial information processing and mapping by R outperform those existing in the traditional software, due to the automation of data processing and computations which is achieved by the functionality of this high-level language. The approaches of map algebra utilise scripts to process remotely sensed data using interactive steps that are completely different from the built-in functions of the traditional software since they use the syntax of R rather then GUI interface. 

Scripting method of the satellite image processing is presented in this paper with an case study of three cities located along the Congo River, central Africa. We have presented a method based on R programming language which offers an integration of its several libraries and applies algorithms for band arithmetics, classification and computing the NDVI indices using unsupervised clustering. A special focus is given to monitoring vegetation and tropical rainforests of Congo. Empirically, we showed that our approach can determine the complexity of landscapes through the classification of pixels and grouping them into 10 land cover classes using the selected functions of RSToolbox as operators for band calculations. Moreover, we performed the experiments on cartographic data processing by R using a series of Landsat-8 OLI/TIRS images and compared landscape dynamics. We observed land cover changes since 2013 until 2022, which illustrate the deforestation in tropical regions of Congo, as reflected in the NDVI for the three target locations: Bumba, Basoko and Kisangani.

The proposed method of using programming for satellite image classification is not a replacement for the state-of-the-art methods of remote sensing data processing. The existing solutions based on GIS methods can be used a quality and conceptually sound methods of image processing in remote sensing, given proper care in supervised classification. A variety of the efficient cases on Landsat images processing exist as reported in many studies. Importantly, however, the automation of image classification is still a challenging problem due to presence of complex background noise on the image scene and restricted functionality of the conventional tools in many GIS. In contrast to the traditional software, R support deep learning analysis through several packages which enhance data processing and enables semantic segmentation. The examples include 'deepnet' toolkit and its dependent library 'deepr', which enhance the training or predicting process or fine-tuning machine learning methods by 'RcppDL'. As such, using deep learning by R can be used with the aim at simplifying the workflow of the satellite image pre- and post-processing, optimising the operations and abstraction of objects during classification, which are required by a variety of remote sensing applications.

In this paper, we have proposed an automatic image processing using R scripts for image classification based on the k-means clustering algorithm in RSToolbox library and band algebra for NDVI computations by terra package. We considered the algorithms of R by computing the raster data and visualising changes in current forest stage on the built maps. The experimental results on vegetation index for the three different datasets taken on years 2013 and 2022 use attributes of various land cover classes for visualising the changes in the land cover types of forest regions of central Congo. A scripting approach to image clustering show that the technique of R programming is able to divide the image scene on the defined number of classes using the unsupervised \emph{k}-means algorithm of RSToolbox. We demonstrated a set of images based on the high-level classification of the Landsat-8 OLI-TIRS raster imagery. Based on the analysis of these images, we confirm that the distribution of the vegetation in rainforest regions of central Africa experience changes. The driving factors are diverse and include both social and climatic processes, which affect tropical forests and contribute to deforestation. The images were processed automatically and compared for the target regions situated along the middle part of Congo River Basin. The results of R performance demonstrated the probabilistic model of vegetation distribution over the study areas based on robust image segmentation. 

The libraries of R language offer improved performance on image processing through a powerful framework for raster data processing. In particular, R presents practical applications for modelling real-time Earth observation data from space borne satellite missions. In particular, we have explored the suitability of R programming for processing Landsat-8 OLI-TIRS images which contain 11 bands which correspond to various wavelengths and 30-m resolution. Unlike the conventional GIS, limited to the default software parameters of the embedded graphical user interface, R proposes flexible command-line solutions. These include extended operational support of bitmap image processing, diverse suitability of map algebra and computations of bands, as shown in this paper. The limitations of the presented work include the restricted approach of geometric and spectral characteristics of bands. Future works can be extended towards the semantic segmentation of the detected land cover classes using deep learning methods, as well as the analysis of the topology of classes for object based image analysis.

%%%%%%%%%%%%%%%%%%%%%%%%%%%%%%%%%%%%%%%%%%
%\vspace{6pt} 

\authorcontributions{Software, methodology, data curation, formal analysis, validation, visualisation---P.L.; 
Supervision, project administration, conceptualisation, resources, investigation, funding acquisition---O.D.; writing---original draft preparation, P.L.; writing---review and editing, P.L. and O.D. All authors have read and agreed to the published version of the~manuscript.}

\funding{This research was funded by the Editorial Office of the Applied Sciences, Multidisciplinary Digital Publishing Institute (MDPI), by providing 100\% discount for APC for publication of this manuscript.}

\institutionalreview{Not applicable}

\informedconsent{Not applicable}

\dataavailability{Not applicable} 

\acknowledgments{We thank the three anonymous reviewers for their careful reading, critical reviews and constructive comments that helped to improve the previous version of the manuscript.}

\conflictsofinterest{The authors declare no conflict of interest} 

\abbreviations{Abbreviations}{
The following abbreviations are used in this manuscript:\\

\noindent 
\begin{tabular}{@{}ll}
CPT & Colour Palette Table \\
DCW & Digital Chart of the World \\
DEM & Digital Elevation Model \\
DN & Digital Number \\
GCP & Ground Control Points \\
GEBCO & General Bathymetric Chart of the Oceans \\
GDAL & Geospatial Data Abstraction Library \\
GIS & Geographic Information System \\
GLS & Global Land Survey \\
GloVis & Global Visualization Viewer \\
GMT & Generic Mapping Tools \\
GRFM & Global Rain Forest Mapping \\
Landsat 8-9 OLI/TIRS & Landsat 8-9 Operational Land Imager and Thermal Infrared \\
LAI & Leaf Area Index \\
LiDAR & Light Detection and Ranging \\ 
L1TP & {Level 1 Terrain Precision (Corrected)} \\
NASA & National Aeronautics and Space Administration \\
NDVI & Normalized Difference Vegetation Index \\
NGA & National Geospatial-Intelligence Agency \\
NIR & Near Infrared \\
REDD & Reducing Emissions from Deforestation and Forest Degradation \\
RMSE & Root Mean Square Error \\
SAR & Synthetic Aperture Radar \\
SRTM & Shuttle Radar Topography Mission \\
SWIR & Short Wave InfraRed \\
UAV & Unmanned Aerial Vehicle \\
USGS & United States Geological Survey \\
UTM & Universal Transverse Mercator \\
WRS & Worldwide Reference System \\
\end{tabular}
}

\begin{adjustwidth}{-\extralength}{0cm}
\reftitle{References}

%=====================================
% References, variant A: external bibliography
%=====================================
\bibliography{Congo.bib}
%\nocite{*}

\section*{Short Biography of Authors}
\bio
{\raisebox{-0.35cm}{\includegraphics[width=3.5cm,height=5.3cm,clip,keepaspectratio]{Photo_author_1.jpg}}}
{\textbf{Polina Lemenkova} is an Engineer and a PhD student of Prof. Dr. O. Debeir in the Image Analysis Research Unit (LISA-IA) in École polytechnique de Bruxelles, Université Libre de Bruxelles. Her research interests are focused on image processing, cartography, spatial data analysis, remote sensing and applied programming. She has a strong research experience in Earth and environmental sciences, mapping and data visualization, specifically in geological and geophysical domains. Her current topic is on seismic data processing using machine learning methods and scripting languages.}
%
\bio
{\raisebox{-0.35cm}{\includegraphics[width=3.5cm,height=5.3cm,clip,keepaspectratio]{Photo_author_2.jpg}}}
{\textbf{Olivier Debeir} is a Head of the Image Analysis Research Unit (LISA-IA) in the École polytechnique de Bruxelles (Brussels Faculty of Engineering), Université Libre de Bruxelles. His research interest is in addressing questions related to artificial intelligence and machine learning in general, and specifically on computer vision, graphics, image analysis, segmentation and classification, pattern recognition and deep learning. His research activity also concerns applied statistics and diverse algorithms of data analysis and modelling in natural sciences. He is an expert in programming languages.}

\end{adjustwidth}
\end{document}

